\documentclass[a4paper, 12pt]{article}
\usepackage[a4paper,
 total={170mm,257mm},
 left=20mm,
 top=20mm]{geometry}
\usepackage[utf8]{inputenc}
\usepackage{amsfonts, amsmath, amsthm, graphicx, lipsum}
\usepackage{fontawesome5}

\usepackage[backend=biber,bibstyle=ieee,dashed=false,sorting=ydnt,sortcites=true,natbib=true,defernumbers=true,maxbibnames=99,giveninits=true,uniquename=false]{biblatex}
\bibliography{pub}

\usepackage{fancyhdr}

\usepackage{hyperref}
\hypersetup{
    colorlinks=true,
    linkcolor=blue,
    urlcolor=blue,
    pdftitle={Awnon CV}
}

\pagestyle{fancy}
\fancyhf{}
\headheight=40pt
\lhead{Curricular Vitae}
\rhead{Awnon Bhowmik}
\chead{\thepage}

\thispagestyle{empty}

% \renewcommand{\bibliography}{heading=none}
\renewcommand*{\mkbibnamegiven}[1]{
\ifitemannotation{highlight}
{\textbf{#1}}
{#1}}
\renewcommand*{\mkbibnamefamily}[1]{
\ifitemannotation{highlight}
  {\textbf{#1}}
  {#1}%
\ifitemannotation{first}
  {\textsuperscript{*}}
  {}%
\ifitemannotation{corresponding}
  {$^\dagger$}
  {}%
}%%

\begin{document}

% \defbibheading{bibliography}{\section*{}}
% \bibliographystyle{ieeetr}

\begin{flushleft}
\large{\textbf{AWNON BHOWMIK}}
\end{flushleft}

\begin{flushleft}
\rule{\textwidth}{1pt}
\textbf{CONTACT INFORMATION}\\
\quad\faMapMarker*\quad University of Central Florida\\
\quad\qquad Department of Mathematics, College of Sciences\\
\quad\qquad 4393 Andromeda Loop N, Orlando, FL 32816\\
\quad\qquad Room 418\\~\\
\quad\faPhone*\quad+1 929 462 8832\qquad\qquad\qquad\qquad\faGithub\quad\href{https://github.com/awnonbhowmik}{awnonbhowmik}\\
\qquad\quad +1 347 287 7680\qquad\qquad\qquad\qquad\faLinkedin\quad\href{https://www.linkedin.com/in/awnon-bhowmik/}{awnon-bhowmik}\\~\\
\quad\faEnvelope\quad\href{mailto:awnon.bhowmik@ucf.edu}{awnon.bhowmik@ucf.edu}\qquad\qquad\,\,\faGoogle\quad\href{https://scholar.google.com/citations?user=nEdZAFkAAAAJ&hl=en}{Awnon Bhowmik}\\
\quad\qquad\href{mailto:abhowmik901@york.cuny.edu}{abhowmik901@york.cuny.edu}\qquad\,\,\,\,\faOrcid\quad\href{https://orcid.org/0000-0001-5858-5417}{Awnon Bhowmik}\\~\\
\quad\faGlobe\quad\href{www.awnonbhowmik.com}{www.awnonbhowmik.com}\\
\rule{\textwidth}{1pt}
\end{flushleft}

\begin{flushleft}
\faLightbulb\quad\textbf{RESEARCH INTERESTS}
\begin{table}[h!]
    \centering
    \begin{tabular}{llll}
      Number Theory & Cryptography & Fractal Geometry\\
      Computational Mathematics & Mathematical Modeling & Algorithm Design \& Analysis
    \end{tabular}
    % \caption{Caption}
    % \label{tab:my_label}
\end{table}
\end{flushleft}

\begin{flushleft}
\faGraduationCap\quad\textbf{EDUCATION}\\
\begin{table}[h!]
    \centering
    \begin{tabular}{llll}
         \textbf{Degree} & \textbf{Institution} & \textbf{Concentration} & \textbf{Obtained}\\
    Ph.D. & University of Central Florida & Mathematics & Expected $2026$\\
    MS & The City College of New York & Mathematics & December $2021$\\
    BS & CUNY York College & Major: Mathematics & January $2020$\\
    & & Minor: Computer Science &\\
    AS & CUNY BMCC & Mathematics & August 2015
    \end{tabular}
    % \caption{Caption}
    % \label{tab:my_label}
\end{table}
\end{flushleft}

\begin{flushleft}
\faBookOpen\quad\textbf{PUBLICATIONS}
\nocite{*}
\printbibliography[heading=none]
\end{flushleft}

\begin{flushleft}
\faAtom\quad\textbf{RESEARCH EXPERIENCE}
\begin{flushleft}
\textbf{Analysis of Searching \& Sorting Algorithms}\hfill November $2019 - $December $2019$
\begin{itemize}
    \item Using files with a varying number of inputs, applied linear and binary search. Analyzed the
execution time for different inputs and proved algebraically and graphically $-$ the time
complexity of those algorithms.
    \item Implemented various common sorting algorithms such as bubble sort, merge sort, etc. on
files containing varying amounts of data. Analyzed the execution time for different inputs
and proved algebraically and graphically $-$ the time complexity of those algorithms.
\end{itemize}
\textbf{Advisor: Dr. Louis D` Alotto}\\
Department of Mathematics \& Computer Science, CUNY York College
\end{flushleft}

\begin{flushleft}
\textbf{Battery Powered Car Problem}\hfill September $2019$
\begin{itemize}
    \item Given a pile of $n$ batteries assuming an ideal mileage of $k$ kilometers per battery, derived a recursive function that models the effective distance traveled by car since it can only carry one spare battery in it along with the other that powers its engine.
    \item Came up with two different algorithms and implemented using C++ code, to demonstrate
    the validity of the mathematical perspective. Illustrated with graphs to further support the claims.
\end{itemize}
\textbf{Advisor: Dr. Louis D` Alotto}\\
Department of Mathematics \& Computer Science, CUNY York College
\end{flushleft}

\begin{flushleft}
\textbf{CSTEP Research Project}\hfill September $2014- $August $2015$\\
\textbf{The Electrolytic Properties of Protein in the Cell Membrane Channel}
\begin{itemize}
    \item Studied the Gramicidin A ion channel of the bacillus Brevis bacteria cell, containing the major amino acid Tryptophan.
    \item Modeled and conjectured that the conformational change in the Tryptophan produces changes in the cell membrane, thus behaving like a capacitive dielectric thereby influencing the ion current.
    \item Applied classical physics to a biological problem to mathematically model the ion current under the influence of a variable membrane dielectric, as a function of amino acid conformation.
\end{itemize}
\textbf{Advisor: Dr. Brett A. Sims}\\
Department of Mathematics, CUNY BMCC
\end{flushleft}

\newpage
\begin{flushleft}
\textbf{Research Project: Safe Distance}\hfill November $2014 - $December $2014$
\begin{itemize}
    \item Modeled the relationship between the initial speed of two cars, the reaction time, and braking distance, relating to "average car lengths" as the unit of measure, essentially proving the "Two Second Rule" implemented by the New York City Department of Motor Vehicles (NY DMV).
    \item Used classical probability and graphs to determine that this rule may change during weathers that hinder visibility and that the distance between two cars must be greater in monsoon and winter than in summer $-$ before the brakes are applied to avoid a collision.
\end{itemize}
\textbf{Advisor: Dr. David Waldman}\\
Department of Science, CUNY BMCC
\end{flushleft}
\end{flushleft}

\begin{flushleft}
\faSuitcase\quad\textbf{WORK EXPERIENCE}
\end{flushleft}

\begin{flushleft}
\faUser\quad\textbf{AFFILIATIONS}
\begin{itemize}
    \item American Mathematical Society
    \item Society of Industrial \& Applied Mathematics
    \item Phi Theta Kappa
\end{itemize}
\end{flushleft}

\begin{flushleft}
\faTrophy\quad\textbf{ACHIEVEMENTS}
\begin{itemize}
    \item Dean’s List, Fall $2019\,-\,$CUNY York College
    \item Dean’s List, Winter $2015$, Spring $2015$, Fall $2014\,-\,$CUNY BMCC
    \item Inaugural Inter\,$-$\,University Math Olympiad Champion, University of Dhaka, July 2011
\end{itemize}
\end{flushleft}
\end{document}